% Root splice: replaces Statement C8, not the surrounding subsections.
\paragraph{Statement C8 (Proved: boundary nonvanishing in every dimension;
local injectivity refuted).}\label{res:C8}
Let $U/K$ be an elementary anabelian variety, namely a finite tower of
relative hyperbolic curves over a field of characteristic zero, and let
$U\hookrightarrow X$ be a smooth proper compactification with
normal-crossings boundary. For a rational boundary point $x$, write
$I(x)$ for the nonempty set of geometric local branches at $x$, and put
$T=\widehat{\mathbf Z}(1)$. The geometric group of the punctured strict
henselian neighbourhood is $T^{I(x)}$, and its natural map
\[
 j_x:T^{I(x)}\longrightarrow\pi_1(U_{\overline K})
\]
has nonzero image. More precisely, its restriction to the diagonal
$d_x:T\to T^{I(x)}$ has nonzero image. Every section through the
punctured henselian neighbourhood therefore has nontrivial cyclotomic
centralizer. When all local branches are rational, that centralizer
contains the whole image of $j_x$; in general it contains the nonzero
image of $j_xd_x$. These assertions hold over every finite base extension
and hence for germs. The assertion that $j_x$ is injective for every
such compactification is false, already on a blown-up product surface.
\passage{Section 4, manuscript pp.~63--65, scans 128--131, especially
the withdrawn assertion $1^\circ$ and the margin at scan 130.}
\begin{proof}
The punctured henselian neighbourhood is
$U_{(x)}=U\times_X\operatorname{Spec}\mathcal O^h_{X,x}$.
Its local exact sequence and action are computed in
Carlson--Corwin~\cite[Proposition 3.4]{B_CarlsonCorwin2026}:
the geometric kernel is
\[
 \operatorname{Hom}\bigl(\mathbf Z[I(x)],T\bigr)=T^{I(x)},
\]
with Galois action obtained by permuting the branches and by the
cyclotomic action on $T$. The natural morphism $U_{(x)}\to U$ supplies
$j_x$ and makes it equivariant for every chosen local section. If the
branches are rational, the cyclotomic kernel fixes this entire local
geometric group. Without that restriction it still fixes the diagonal,
because every permutation fixes a constant tuple. It remains to prove
that the diagonal survives the displayed map.

We use the following precise valuative input. Hoshi
\cite[Lemma 2.8, pp.~175--176]{B_Hoshi2014} treats a relative hyperbolic
curve $Y\to S$, where $S$ is connected, noetherian, separated and
normal over a characteristic-zero field. For a strictly henselian
discrete valuation ring $R$ over $S$, with fraction field $F$, a given
$S$-map $\operatorname{Spec}F\to Y$ extends over
$\operatorname{Spec}R$ exactly when its homomorphism
$\pi_1(\operatorname{Spec}F)\to\pi_1(Y)$ has trivial image.

This gives the same extension criterion for a hyperbolic polycurve over
an algebraically closed characteristic-zero field. Here is the
induction, including the maps involved. For a tower
$U=U_n\to U_{n-1}\to\cdots\to U_0=\operatorname{Spec}\overline K$
and $a:\operatorname{Spec}F\to U$ with trivial fundamental-group
image, its composite to $U_{n-1}$ also has trivial image. By induction
that composite extends to $b:\operatorname{Spec}R\to U_{n-1}$.
Regard $R$ as a ring over the normal base $U_{n-1}$ by $b$.
Hoshi's lemma applied to the last relative curve and to the original
map $a$ extends $a$ to $U_n$. The base case is a point. The converse
follows from $\pi_1(\operatorname{Spec}R)=1$, and uniqueness follows
from separatedness.

Apply this criterion to an explicit arc at the chosen geometric point
$\bar x$ over $x$. Let $r=|I(x)|$ and $d=\dim U$. Choose branch
parameters $z_1,\ldots,z_r$ in the strict henselian local ring and
complete them to a regular system of parameters. Normal crossings and
smoothness identify its completion with
\[
 \widehat{\mathcal O}_{X_{\overline K},\bar x}
   \simeq\overline K[[z_1,\ldots,z_d]],
 \qquad X\setminus U:\ z_1\cdots z_r=0.
\]
The continuous local homomorphism to $R=\overline K[[t]]$ defined by
\[
 z_i\longmapsto t\quad(1\le i\le r),\qquad
 z_i\longmapsto0\quad(r<i\le d)
\]
defines an arc $a:\operatorname{Spec}R\to X_{\overline K}$ through
$\bar x$. Its generic point lands in $U_{\overline K}$ because the
boundary equation pulls back to the nonzero element $t^r$.
Since $R$ is strictly henselian, the arc factors through the strict
henselization at $\bar x$. On punctured neighbourhoods the covers
$z_i^{1/n}$ pull back to $t^{1/n}$ for every $i$ and every $n$.
Consequently the induced geometric homomorphism is precisely
\[
 \pi_1(\operatorname{Spec}\overline K((t)))
   =T\xrightarrow{d_x}T^{I(x)}\xrightarrow{j_x}\pi_1(U_{\overline K}).
\]
If $j_xd_x$ were zero, the extension criterion just proved would give
$\widetilde a:\operatorname{Spec}R\to U_{\overline K}$ with the same
generic-point map as $a$. Its composite with $U_{\overline K}\to
X_{\overline K}$ equals $a$: two maps from this integral trait to the
separated scheme $X_{\overline K}$ which agree at the generic point
agree everywhere. Their closed point would then lie both in
$U_{\overline K}$ and at $\bar x\in X_{\overline K}\setminus
U_{\overline K}$, a contradiction. Thus $\operatorname{im}(j_xd_x)$
is nonzero. The previously proved equivariance now gives the asserted
nontrivial subgroup of the cyclotomic centralizer.

For the failure of injectivity, take
$U=(\mathbf P^1\setminus\{0,1,\infty\})^2$, compactify in
$(\mathbf P^1)^2$, and blow up $(0,2)$. With old local coordinates
$t$ and $u=z_2-2$, the old boundary is $t=0$. In the chart
$t=ab$, $u=b$, the new boundary has the two branches $a=0$ and $b=0$.
Each branch loop maps to the loop around $t=0$, while the loop around
$u=0$ contracts in the original neighbourhood. Thus the actual map is
\[
 T^2\longrightarrow I_0\subset\pi_1(U_{1,\overline K}),
 \qquad(m,n)\longmapsto m+n.
\]
Its kernel is $\{(m,-m):m\in T\}$ and its image is $I_0\ne1$.
The diagonal maps by multiplication by $2$, so it is indeed nonzero.
The finite matrix reductions in \texttt{code/root\_math\_checks.py}
record this kernel and image modulo $2,3,5,25$.
\end{proof}

The converse requires a geometric boundary centre and a factorization
through its actual punctured henselian group for every section with
nonzero cyclotomic centralizer, with a specified compactification and
arithmetic domain. Merely recognizing a nontrivial abelian subgroup
does not construct that factorization. This obligation already contains
the curve section-surjectivity problem through explicit product maps.
Indeed, let $C/K$ be any hyperbolic curve, let $Q=\mathbf
P^1\setminus\{0,1,\infty\}$, and fix a rational cusp section $t$ of
$Q$. For every section $\sigma$ of $C$, the product section
$(\sigma,t)$ of $C\times Q$ has a centralizer containing
$\{1\}\times I_0$. A boundary-local factorization of that section in
the product compactification would, by the first projection, express
$\sigma$ either as the point section at an interior first coordinate
or as a cuspidal section at a boundary first coordinate. Thus a
universal boundary-factorization assertion for these products entails
the point-or-cusp assertion for all such $\sigma$ over the same field.
For a statement about germs, the same argument retains the corresponding
finite-extension domain. No such exhaustion of the section set is
asserted here.

For source comparison, Carlson--Corwin
\cite[Theorem 3.7 and Proposition 4.12]{B_CarlsonCorwin2026} prove
compactification independence of the \emph{closure} of the boundary
section set and an inclusion into the \emph{closure} of the inductively
defined cuspidal set of a tower. The closure bars are essential in both
statements. Their Proposition 3.4 supplies the precise local group and
action used above; the nonvanishing argument here additionally uses
Hoshi's valuative lemma. None of these comparisons supplies the
universal factorization in the preceding paragraph.
